% Lecture 03: Указатели, память и массивы

\section{Лекция 3. Указатели, память и массивы}
\label{sec:lecture-03}

В предыдущей лекции объект хранил одно значение, а цикл повторял действие.
Теперь нужно хранить несколько однотипных объектов подряд и понимать их
адреса. Сначала построим честную модель raw array и указателя, затем увидим,
какую ручную работу снимают \texttt{std::array} и \texttt{std::vector}.

\begin{importantbox}[Главный вопрос]
Что означает «несколько объектов подряд в памяти», почему возникает pointer
arithmetic и почему для обычного изменяемого набора данных современный C++
чаще выбирает \texttt{std::vector}?
\end{importantbox}

\subsection{Объект, storage, lifetime и адрес}

Уточним три разных понятия: \emph{storage} --- память, достаточно большая и
правильно выровненная для объекта; \emph{object lifetime} --- интервал
выполнения, в котором объект существует; \emph{value} --- информация, которую
он хранит. Storage само по себе ещё не всегда означает живой объект. Для
обычного локального объекта lifetime начинается после инициализации и
заканчивается при выходе из блока.

\begin{cppcode}
int score{42};
std::cout << score << '\n';   // значение
std::cout << &score << '\n';  // адрес объекта
\end{cppcode}

\begin{definition}
\textbf{Адрес} --- значение, обозначающее положение объекта или функции в
абстрактной модели C++. Указатель может хранить такое значение.
\end{definition}

Выведенный адрес помогает наблюдать взаимное размещение, но его нельзя без
оговорки называть физическим адресом RAM. Automatic storage duration также не
обязывает реализацию держать объект на CPU stack: это обычная реализация, а
не требование стандарта.

\texttt{sizeof(expression)} даёт размер представления в байтах типа
\texttt{std::size\_t}. Байт C++ равен размеру \texttt{char} и содержит не
менее восьми бит; универсального обещания ровно восьми бит нет. Обычно операнд
\texttt{sizeof} не вычисляется. Результат «\texttt{int} равен 4 байтам» можно
наблюдать на машине, но нельзя превращать в закон языка. Alignment означает,
что адрес некоторых типов должен быть подходяще выровнен; компилятор учитывает
это при размещении, подробности сейчас не нужны.

\subsection{Pointer basics: адрес и значение по адресу}

\begin{cppcode}
int score{42};
int* pointer{&score};

std::cout << pointer << '\n';   // pointer value: адрес score
std::cout << *pointer << '\n';  // pointed-to value: 42
*pointer = 50;                  // меняет score
\end{cppcode}

\texttt{int* pointer} объявляет объект типа «pointer to \texttt{int}».
Звёздочка в declaration участвует в записи типа; \texttt{*pointer} в
expression выполняет \emph{dereference}. Оператор \texttt{\&score} берёт
адрес. Pointer --- не «просто integer»: у него отдельный тип, допустимые
операции и правила действительности.

\begin{whybox}[Зачем различать pointer и pointed-to value]
Копирование \texttt{pointer} копирует адресное значение, не \texttt{score}.
Присваивание \texttt{*pointer = 50} работает с объектом по адресу.
\end{whybox}

Минимальный \texttt{const} для чтения без изменения:

\begin{cppcode}
int score{42};
const int* view{&score};
// *view = 50; // compile-time error
\end{cppcode}

Сам \texttt{view} можно направить на другой \texttt{int}. Полная система
const будет в следующем блоке.

\subsubsection{nullptr и действительность}

\texttt{nullptr} явно означает «сейчас не указывает на объект» и имеет
специальный тип. Старый литерал \texttt{0} не используем как основной стиль.

\begin{cppcode}
int* selected{nullptr};
if (selected != nullptr) {
    std::cout << *selected << '\n';
}
\end{cppcode}

Dereference требует подходящий живой объект. Проверка на \texttt{nullptr}
устраняет только null case, но не распознаёт dangling pointer.

\begin{definition}
\textbf{Dangling pointer} сохраняет адресное значение после окончания lifetime
объекта, на который раньше указывал. Доступ к прежнему объекту через него
недопустим.
\end{definition}

Неинициализированный локальный pointer фундаментального типа имеет
неопределённое значение; его использование может иметь undefined behavior.
Безопасная привычка --- сразу дать адрес живого объекта или \texttt{nullptr}.

\subsection{Raw array: фиксированная последовательность}

\begin{cppcode}
int temperatures[4]{18, 20, 21, 19};
\end{cppcode}

Это built-in, или raw, array из четырёх \texttt{int}. Extent \texttt{4}
входит в array type и известен при компиляции. Все элементы имеют один
тип, индексы идут от нуля, объекты расположены непрерывно без промежутков
между соседними элементами.

\begin{cppcode}
int values[4]{};          // все элементы равны нулю
int inferred[]{2, 4, 6}; // extent 3 выведен из initializer
int partial[4]{7, 8};    // 7, 8, 0, 0
values[0] = 10;
\end{cppcode}

Во время выполнения \texttt{values[i]} выбирает элемент, но built-in access
не вставляет обязательную проверку границ.

\begin{cppcode}
int values[]{10, 20, 30, 40};
const std::size_t count = sizeof(values) / sizeof(values[0]);
\end{cppcode}

\texttt{sizeof(values)} здесь означает весь array object, а
\texttt{sizeof(values[0])} --- элемент. Деление --- лишь локальный
учебным приёмом там, где expression всё ещё имеет array type.

\begin{ubbox}[Out-of-bounds]
Для extent $N$ действительны индексы от 0 до $N-1$. \texttt{values[N]} уже
за границей. Чтение или запись имеет UB. UB не обещает crash: возможны
правдоподобный вывод, повреждение состояния или иное поведение после
оптимизации.
\end{ubbox}

\subsection{Array-to-pointer conversion: массив не pointer}

\begin{cppcode}
int values[]{10, 20, 30, 40}; // type: int[4]
int* first{values};           // conversion; type: int*
\end{cppcode}

Array object и pointer object различны. У массива четыре элемента и fixed
extent; у \texttt{first} одно pointer value. Во многих expressions array
неявно преобразуется в pointer на первый элемент: при инициализации
\texttt{first}, в \texttt{values + 1} и при передаче обычному pointer
parameter. Функции подробно появятся в блоке~4.

Преобразование происходит не везде:

\begin{cppcode}
std::cout << sizeof(values) << '\n'; // весь int[4]
std::cout << sizeof(first) << '\n';  // один int*
std::cout << &values << '\n';        // pointer на whole array, другой тип
\end{cppcode}

После conversion pointer не несёт extent. Поэтому деление размера pointer на
размер элемента не находит длину исходного массива. Pointer можно
направить заново, имя массива --- нельзя; raw arrays также не присваиваются
целиком. Фраза «имя массива --- pointer» скрывает эти различия. Точно:
array expression во многих контекстах \emph{преобразуется} в pointer.

\subsection{Pointer arithmetic существует из-за массивов}

\begin{cppcode}
int values[]{10, 20, 30, 40};
int* p{values};
std::cout << *(p + 2) << '\n';
++p;
std::cout << p - values << '\n';
\end{cppcode}

Для \texttt{int*} операция \texttt{p + 1} означает следующий \texttt{int},
не следующий машинный байт. \texttt{a[i]} имеет семантику
\texttt{*(a + i)}. Вычитание pointers даёт число элементов типа
\texttt{std::ptrdiff\_t}, но оба должны относиться к одному array object.

\begin{cppcode}
int values[]{10, 20, 30, 40};
int* begin{values};
int* end{values + 4};
for (int* current = begin; current != end; ++current) {
    std::cout << *current << '\n';
}
\end{cppcode}

\texttt{end} --- one-past pointer: его разрешено сформировать и сравнить как
границу полуинтервала $[begin,end)$, но \texttt{*end} имеет UB. Нельзя уходить
дальше one-past или выполнять арифметику между несвязанными arrays.

\begin{performancebox}[Индекс или pointer increment?]
\texttt{a[i]} и \texttt{*(a + i)} выражают ту же адресацию. Нельзя обещать,
что одна запись быстрее: нужен benchmark или assembly для конкретных compiler,
flags и workload. В обычном коде выбираем ясность.
\end{performancebox}

\subsection{Коротко о C-style strings}

\begin{cppcode}
char label[]{'C', '+', '+', '\0'};
char same_label[]{"C++"};
\end{cppcode}

C-style string --- последовательность \texttt{char} с null terminator
\texttt{'\textbackslash 0'}. Extent второго массива равен 4, хотя видимых
символов три. Отсутствие terminator позволяет C API читать за границу. Для
совместимости такая форма важна, но основной текстовый тип курса ---
\texttt{std::string}.

\subsection{std::array: удобный fixed-size container}

\begin{cppcode}
#include <array>
std::array<int, 4> values{10, 20, 30, 40};
std::cout << values.size() << '\n';
std::cout << values[1] << ' ' << values.at(2) << '\n';
std::cout << values.data() << '\n';
\end{cppcode}

Размер входит в тип и не меняется runtime, storage contiguous.
Объект знает \texttt{size()}, копируется и присваивается целиком;
\texttt{at()} проверяет границу и при ошибке сообщает exception. Исключения
будут отдельной темой. \texttt{data()} даёт pointer на первый элемент, а
\texttt{begin()}/\texttt{end()} интуитивно обозначают начало и one-past;
iterators системно изучим позже.

\subsection{std::vector: dynamic size и владение storage}

Raw array не растёт. Ручной dynamic array потребовал бы получить новый
участок, перенести элементы и освободить старый на всех путях. Naked
\texttt{new[]} / \texttt{delete[]} сейчас не используем: управление
динамической памятью будет позже.

\begin{cppcode}
std::vector<int> values{10, 20, 30};
values.push_back(40);
std::cout << values.size() << ' ' << values.capacity() << '\n';
std::cout << values[0] << ' ' << values.at(3) << '\n';
\end{cppcode}

Vector владеет внутренним storage и управляет lifetime элементов. Элементы
contiguous, \texttt{data()} возвращает pointer на первый. \texttt{empty()}
проверяет нулевой size. \texttt{operator[]} требует правильной границы от
программиста; \texttt{at()} проверяет её.

\begin{definition}
\texttt{size()} --- число живых элементов. \texttt{capacity()} --- сколько
поместится в текущем storage без нового выделения. Capacity не меньше size,
но позиции между ними ещё не являются элементами.
\end{definition}

Когда capacity исчерпана, \texttt{push\_back} может сделать
\emph{reallocation}: получить больший contiguous участок, перенести элементы
и отказаться от старого. Точный коэффициент роста и layout не гарантированы.

\begin{cppcode}
std::vector<int> values{10};
int* old_data{values.data()};
const std::size_t old_capacity{values.capacity()};
while (values.capacity() == old_capacity) {
    values.push_back(20);
}
// old_data больше не dereference: reallocation его инвалидировал
\end{cppcode}

Элементы теперь живут в новом storage, поэтому старый адрес не обозначает
прежний живой объект. Reallocation инвалидирует pointers, references и
iterators на элементы. Последние два термина здесь лишь анонс будущих тем.

\begin{cppcode}
std::vector<int> values;
values.reserve(100);
std::cout << values.size();     // 0
std::cout << values.capacity(); // не меньше 100
\end{cppcode}

\texttt{reserve} подготавливает capacity, но не создаёт элементы и не меняет
logical size; это не \texttt{resize}. Он уменьшает число reallocations при
хорошем прогнозе, но сам способен инвалидировать старые pointers.

Один рост с reallocation стоит $O(n)$. Capacity растёт с запасом, поэтому
дорогая операция нужна не при каждом добавлении. \texttt{push\_back} имеет
amortized $O(1)$: постоянную среднюю стоимость на длинной серии, не гарантию
$O(1)$ каждого вызова.

\subsection{Raw array, std::array или std::vector?}

\begin{center}\small
\begin{tabularx}{\textwidth}{@{}lXXX@{}}
\toprule
Свойство & raw array & \texttt{std::array} & \texttt{std::vector} \\
\midrule
Размер & fixed & fixed, часть типа & dynamic \\
Знает размер & extent в array type & \texttt{size()} & \texttt{size()} \\
Checked access & нет & \texttt{at()} & \texttt{at()} \\
Copy/assignment & не целиком & целиком & целиком \\
Storage & contiguous & contiguous & contiguous \\
Pointer access & conversion & \texttt{data()} & \texttt{data()} \\
Обычная роль & ABI/C API, low-level & fixed container & растущая sequence \\
\bottomrule
\end{tabularx}
\end{center}

Raw arrays не запрещены: они раскрывают язык и нужны некоторым ABI,
embedded- и C interfaces. Для обычного владения fixed sequence удобнее
\texttt{std::array}, runtime-размером --- \texttt{std::vector}.

\subsection{Карта memory errors}

\textbf{Uninitialized pointer:} неопределённое значение адреса; \textbf{Null dereference:} \texttt{*nullptr} не обозначает объект и имеет UB; \textbf{Dangling/use-after-lifetime:} адрес сохранился после конца lifetime; \textbf{Out-of-bounds:} индекс или dereference вышел за элементы; \textbf{Invalid arithmetic:} pointer вышел за one-past одного массива или операция смешала несвязанные arrays; \textbf{Vector invalidation:} reallocation закончила lifetime элементов в старом storage, поэтому старый pointer нельзя dereference.

Scope отвечает, где имя доступно в тексте; lifetime --- когда объект
существует runtime. Вынос адреса за scope имени не продлевает lifetime.

\subsection{ASan и UBSan: наблюдаем часть UB}

AddressSanitizer добавляет runtime-проверки для части out-of-bounds,
use-after-free и use-after-scope. UndefinedBehaviorSanitizer проверяет часть
других UB. Это instrumentation toolchain, не изменение стандарта.

\begin{shellcode}
cmake -S . -B build-sanitize -DENABLE_SANITIZER_CASES=ON \
  -DCMAKE_BUILD_TYPE=Debug
cmake --build build-sanitize
./build-sanitize/raw_oob
\end{shellcode}

PASS требует реальных причин: stack-buffer-overflow, stack-use-after-scope и
heap-use-after-free для pointer после vector reallocation. Эти программы
намеренно имеют UB и изолированы от safe lab.

\begin{warningbox}[Sanitizer не доказывает корректность]
Instrumentation видит не все нарушения и только на выполненных путях.
Отсутствие сообщения не доказывает отсутствие UB.
\end{warningbox}

\subsection{Производительность без лозунгов}

Все три варианта дают contiguous элементы. Последовательный проход часто
помогает cache locality, но конкретный cache --- свойство реализации и
железа. Vector хранит metadata; обычно она мала, но стандарт не обещает layout
или число слов. Reallocation стоит $O(n)$, reserve может уменьшить allocations,
а проход по готовому vector может быть столь же cache-friendly, как по array.

Поэтому «vector медленнее массива» не универсально. Нужно назвать operation,
size, access pattern, compiler flags и платформу, затем измерить. Benchmark
здесь N/A: лаборатория наблюдает семантику storage, не выдаёт нестабильные
наносекунды за закон.

\subsection{Эксперимент 1: course\_inventory}

Read-only утилита в
\texttt{examples/03\_pointers\_memory\_and\_arrays/01\_course\_inventory}
проходит по фиксированным 15 slugs в \texttt{std::array}, реально проверяет
lecture file/examples directory, считает байты, строки и обычные файлы.
Измеренные строки собираются в \texttt{std::vector<BlockInfo>}.

Адреса \texttt{blocks.data()} и \texttt{rows.data()} наблюдают начало двух
contiguous областей, но не физический RAM address. Fixed input и dynamic
results показывают выбор контейнера; reserve меняет capacity, не size.

\begin{shellcode}
cd examples/03_pointers_memory_and_arrays/01_course_inventory
cmake -S . -B build
cmake --build build
./build/course_inventory --root ../../..
./build/course_inventory --root ../../.. --block 3
ctest --test-dir build --output-on-failure
\end{shellcode}

PASS: вывод основан на текущих файлах, totals согласованы, один блок и все
блоки завершаются с exit code 0. Утилита ничего не изменяет.

\subsection{Эксперимент 2: array\_vs\_vector\_lab}

Safe executable показывает адреса raw array, оба \texttt{sizeof}, conversion,
pointer arithmetic и one-past без dereference; затем fixed
\texttt{std::array}; затем size/capacity, смену \texttt{data()} при vector
reallocation и стабильность в пределах reserve.

\begin{shellcode}
cd examples/03_pointers_memory_and_arrays/02_array_vs_vector_lab
cmake -S . -B build
cmake --build build
./build/array_vs_vector_lab
ctest --test-dir build --output-on-failure
\end{shellcode}

Safe run не dereference старый pointer. Три отдельные программы намеренно
показывают raw out-of-bounds, dangling access и vector invalidation. PASS ---
успешный safe test и три ожидаемые sanitizer diagnostics с ненулевыми exit
codes, а не простое наличие исходников.

\subsection{Сильная сторона C++ и её цена}

C++ позволяет видеть адреса и стоимость роста, затем выбрать abstraction без
отказа от contiguous storage. Vector снимает ручное ownership, но не отменяет
lifetime/invalidation. Raw pointers не несут ownership guarantee; свобода
означает ответственность за границы. Производительность зависит от access
pattern, allocations и locality, а не лозунга «низкоуровневое быстрее».

\subsection{Проверка понимания}

Нужно уметь объяснить:
1) различие объекта, адреса и значения по адресу; 2) почему pointer не просто integer, а array не pointer; 3) где conversion теряет extent и почему различаются два \texttt{sizeof}; 4) границы pointer arithmetic и роль one-past; 5) отличие scope от lifetime; 6) отличие size/capacity и reserve/resize; 7) причинную связь reallocation и invalidation; 8) выбор raw array, \texttt{std::array}, \texttt{std::vector}; 9) что показывает и чего не доказывает sanitizer.

\subsection{Источники для сверки}

Правила lifetime, arrays, conversions, additive operators и \texttt{sizeof}
сверены с C++17 working draft N4659. Основы сопоставлены с MIT 6.096 Lectures
4, 5 и lifetime-фрагментами Lecture 8, но устаревшие мифы «array is pointer» и
«pointer is just integer» исправлены. Контейнеры сверены с
\textit{Modern C++ Tutorial}, performance context --- с ISO Technical Report
on C++ Performance без превращения implementation facts в гарантии.
